\subsection{All Hopf-leaf stars, including an arbitrary central component}

\begin{statement}{C.2.UN.53}{Proved}{Complete reduced Hopf-star support criterion}
Let $S$ be a specified compact connected SNC surface with smooth
components $P,H_1,\ldots,H_n$, $n\ge1$. Each $H_i$ is a primary
or secondary Hopf surface. Its intersection with $P$ is one
connected smooth curve $J_i$, the $J_i$ are disjoint on $P$,
and there are no other intersections. We impose no initial
classification or elliptic-fibration hypothesis on $P$.

The following two occurrence assertions are equivalent: $S$
is the full reduced central fibre of a proper two-torus
degeneration with smooth total space; and $S$ is such a fibre
of a smooth compact threefold over $\PP^1$, with every other
fibre a smooth two-torus. They hold precisely for the following
explicit data.

Every leaf pair is isomorphic to one of
\[
 \begin{gathered}
 (H^{\rm diag}_{r_i}(\beta_i)/\mu_{n_i},J_i),\qquad
 (H^{\rm nl}_{r_i}(\beta_i)/\mu_{n_i},J_i),\\
 r_i\ge1,\qquad 0<|\beta_i|<1,\qquad \gcd(r_i,n_i)=1 ,
 \end{gathered}                                             \tag{C.AS.1}
\]
where the two contractions on $\C^2\setminus\{0\}$ are,
respectively,
\[
 (z,w)\longmapsto(\beta_i^{r_i}z,\beta_iw),\qquad
 (z,w)\longmapsto(\beta_i^{r_i}z+w^{r_i},\beta_iw),
 \quad
 \zeta(z,w)=(\zeta^{r_i}z,\zeta w),                         \tag{C.AS.2}
\]
and $J_i$ is the image of $w=0$. The normal lines on every
double curve are dual. The central component has a
fixed-Jacobian, identity-integral-monodromy elliptic
quasi-bundle
\[
                         q:P\longrightarrow B             \tag{C.AS.3}
\]
over an elliptic curve. Its marked full fibres are $r_iJ_i$;
there are no other multiple fibres and every fibre support
is smooth. The actual inverse logarithmic transform $T_0$
defined by (C.SN.3)--(C.SN.6), retaining all the specified
axis identifications and the entire analytic torsor,
satisfies
\[
                         c(T_0)=0\in H^2(B,\Lambda_F).     \tag{C.AS.4}
\]
Finally there is a positive integer $R$ such that
\[
 \begin{cases}
 r_i=R,&H_i\text{ is of the nonlinear form in (C.AS.2)},\\
 r_i<R,&H_i\text{ is of the diagonal form in (C.AS.2)}.
 \end{cases}                                             \tag{C.AS.5}
\]
When every leaf is diagonal, the last test is satisfied by
any $R>\max_i r_i$. When a nonlinear leaf occurs, it
determines $R$, and every other nonlinear leaf has that
same resonance.

The conclusion retains the full isomorphism class of $S$.
In particular it retains arbitrary flat torsor parameters,
finite secondary orders, independent primitive angular
directions, and the given elliptic gluing maps. Every support
in the compact realization has Euler characteristic zero.
\end{statement}

\begin{proof}
\par\medskip\noindent\textbf{The canonical identity on an arbitrary Hopf leaf.}
We first use an actual reduced smoothing. Let the divisor
of a local canonical evaluation generator be
$k_PP+\sum_i k_iH_i$. Its existence and the fact that it
has only vertical zeroes are the exact direct-image
calculation in C.2.UN.38. Adjunction on one leaf gives
\[
                  K_H=\OO_H(-aJ),\qquad
                  a=k_H-k_P+1\in\Z.                      \tag{C.AS.6}
\]
We classify this equality on the actual pair $(H,J)$,
including the full secondary deck group.

Take a normal primary cyclic subgroup of the Hopf deck
group $\Gamma$. Its generator $F$ is central: conjugation
can send it only to $F$ or $F^{-1}$, while the nonzero
logarithm of the modulus of its derivative determinant
excludes the latter. This is the exact normal-core
argument of C.2.UN.43. We therefore have a finite
unramified primary cover $H'\to H$, and the full inverse
image $D$ of $J$ is a disjoint union of smooth elliptic
curves. The finite quotient acts transitively on these
curves, since $J$ is connected. Pullback gives
\[
                           K_{H'}=-aD.                  \tag{C.AS.7}
\]

For a nonlinear primary contraction the curve computation
of C.2.UN.38 gives its unique axis, and the full
centralizer computation of C.2.UN.43 gives exactly the
nonlinear quotient in (C.AS.1)--(C.AS.2).
The invariant form $dz\wedge dw/w^{r+1}$ has divisor
$-(r+1)J$, on the quotient as well. Since $\OO_H(J)$ has
infinite order, (C.AS.6) forces $a=r+1$.

Consider a diagonal primary contraction. Its invariant
curve computation can be read directly from an equation
on $\C^2$. Closure across the origin and removal of the
nonconstant unit in its transformation give an eigenfunction
of $F$. Comparison of its power-series monomials gives
the axes in the absence of a multiplicative resonance.
In that case the complete centralizer is diagonal, so
it cannot interchange the axes. Thus $D$ is one axis.
If, for example, it is $w=0$, the exact automorphy
characters in (C.AS.7) give
$\alpha\beta=\beta^a$, which is the resonance
$\alpha=\beta^{a-1}$. The contraction moduli make
$a-1>0$. This contradicts the assumed absence of a
resonance and accounts for that case through an actual
character equality.

In the resonant case pass to a further central cyclic
subgroup to remove the finite root-of-unity discrepancy.
Its generator has the form
\[
                 F(z,w)=(\gamma^p z,\gamma^q w),
       \qquad p,q\ge1,\quad\gcd(p,q)=1,\quad |\gamma|<1 .
                                                               \tag{C.AS.8}
\]
Here the axes have defining weights $p,q$, and every
other irreducible curve has equation $z^q=cw^p$ of
weight $pq$. This follows from the same eigenfunction
calculation: after an axis monomial is removed, a
weighted homogeneous polynomial factors into these
binomials over $\C$. Its irreducible factor is unique
because $p,q$ are coprime. The curves are smooth away
from the origin, as required for the elliptic curves
on $H'$.

The canonical weight is $p+q$. If $D$ is an orbit of
$m$ non-axis curves, (C.AS.7) therefore gives
\[
                              ampq=p+q.                 \tag{C.AS.9}
\]
The same calculation with a single axis replaces $mpq$
by $p$ or $q$. These are genuine character equalities:
a purported extra holomorphic unit and its inverse
extend across the origin, and their nonzero value there
forces the identical character. In particular $a>0$.

For $p,q>1$, coprimality gives $pq>p+q$, so (C.AS.9)
is impossible. A single axis would require $p\mid q$
or $q\mid p$, also impossible. In this case the
centralizer preserves each axis, so an orbit containing
both axes cannot occur. For $p=r>1,q=1$, the curves
$z=cw^r$, including $z=0$, all have weight $r$.
The triangular centralizer can permute these curves;
an orbit of $m$ of them has total weight $mr$ and
would require $amr=r+1$, hence $r\mid1$.
The only possibility is the weight-one axis $w=0$,
and then $a=r+1$.

We retain the full deck group in this last calculation.
Commuting with $(\gamma^rz,\gamma w)$ gives exactly
\[
              (z,w)\longmapsto(Az+Cw^r,Bw).             \tag{C.AS.10}
\]
Its induced action on $z/w^r$ is affine. The central
cyclic subgroup acts trivially there; hence the image
is a finite affine group. Averaging the orbit of zero
gives an actual common fixed point. The corresponding
coordinate change $z\mapsto z-cw^r$ makes every element
diagonal, while retaining the marked axis. Equation
(C.AS.6) then says, for every deck element,
$AB=B^{r+1}$, so $A=B^r$. Thus $\Gamma$ is the
weighted scalar group with an infinite contracting
generator $\beta$ and finite roots of unity
$\mu_n$. Freeness on the marked axis gives
$\gcd(r,n)=1$. This is exactly the diagonal pair in
(C.AS.1)--(C.AS.2).

It remains to check the scalar case $p=q=1$,
where the full secondary group need not initially
be diagonal. Every component of $D$ lifts to a
line, and (C.AS.7) gives $am=2$. If $m=1$,
the finite projective group fixes that line.
On its affine complement, averaging gives a
second fixed point, so the linear group becomes
diagonal in coordinates retaining the first line.
The character equality $AB=B^2$ forces $A=B$
for every element. We obtain the case $r=1$
of (C.AS.2), with scalar finite action.
If $m=2$, then $a=1$ and an element of the full
deck group interchanges the two lines. Choose
them as the two coordinate axes. The unique
log-canonical form up to a constant is
\[
                             {dz\over z}\wedge{dw\over w}.
                                                               \tag{C.AS.11}
\]
Every such interchange multiplies this form
by $-1$. A form with exactly the same divisor
differs by a nowhere-zero holomorphic function
on the compact primary Hopf surface, hence by
a constant. It therefore cannot descend.
This contradicts the actual equality
$K_H+J=0$ in (C.AS.6). Thus this last orbit
does not satisfy (C.AS.6). We have proved,
without an a priori restriction on the secondary
group, precisely the leaf-pair list (C.AS.1),
and in every surviving case
\[
                              k_i-k_P=r_i\ge1.          \tag{C.AS.12}
\]

\par\medskip\noindent\textbf{The central component is forced into the stated class.}
Adjunction on $P$ now gives the exact effective
canonical divisor
\[
                         K_P=\sum_i(r_i-1)J_i .         \tag{C.AS.13}
\]
Thus $\kappa(P)\ge0$ and $p_g(P)\ge1$.
C.2.UN.28 gives $p_g(P)=1$. The surface $P$
is minimal. Indeed, a $(-1)$-curve is rational,
so is distinct from every elliptic $J_i$.
Its intersections with all of them are
nonnegative, contradicting $K_P\cdot E=-1$
and (C.AS.13).

We now specify the covering comparison needed
to apply C.2.UN.31 to this original component.
Its full-fibre multiplicity is one. In the
normalized root base change $t=s^N$, a generic
chart is $x=s^N$ after a holomorphic unit is
absorbed. There is exactly one strict covering
component and its map to $P$ has degree one.
Equivalently, the total generic degree of the
strict component covers is the original
multiplicity, as calculated in C.2.UN.36.
After resolution it remains birational to $P$.
Both minimal models are therefore the same
smooth surface. The actual root-cover conclusion
of C.2.UN.31 thus applies to $P$ itself:
it is a two-torus, or a fixed-Jacobian,
identity-monodromy elliptic quasi-bundle over
a curve of genus at most one, with only smooth
multiple fibres. In the latter description
$p_g(P)=g(B)$, so the base is elliptic.
No assertion about an arbitrary quotient
component has been substituted for this
degree-one comparison.

We also verify that all the given curves
$J_i$ are fibres of one such fibration.
On a torus, an embedded elliptic curve is a
translate of an elliptic subgroup: lift its
map to universal covers and differentiate;
the derivative is periodic and holomorphic,
hence constant. Two elliptic subgroups with
distinct tangent lines give a surjective
isogeny from their product to the two-torus.
Every two translates then meet. The disjoint
$J_i$ consequently have one common tangent
line and are fibres of the actual elliptic
quotient.

In the quasi-bundle description over an
elliptic base the canonical divisor is
$\sum(m-1)F_{\rm red}$, by the exact root
chart calculation in C.2.UN.49.
If a multiple fibre occurs, a horizontal
$J_i$ would have positive intersection with
this canonical divisor. But its torsion
normal line gives $J_i^2=0$, and adjunction
gives $K_P\cdot J_i=0$. Thus every $J_i$
is vertical. If no multiple fibre occurs,
we have a principal elliptic bundle.
A nonzero integral Chern vector excludes
a horizontal elliptic curve: its finite
unramified map to the base would pull the
bundle back to one with a section, while
its Chern vector would be multiplied by
the nonzero covering degree. If that Chern
vector is zero, the exact Cech lattice
calculation (C.NS.14) identifies $P$ with
a torus, reducing to the preceding argument.
This proves (C.AS.3) with every marked
curve retained.

The normal of the leaf axis has exact
torsion order $r_i$, including its finite
secondary character by (C.SN.3)--(C.SN.4).
Normal duality in the reduced smoothing
therefore gives full marked fibre $r_iJ_i$
on $P$. Subtracting the marked terms of
its canonical root-chart formula from
(C.AS.13) leaves an effective sum of
unmarked multiple fibres linearly
equivalent to zero. Such a nonzero
effective divisor on a compact connected
complex surface cannot be principal:
its defining meromorphic function would
have no poles and hence be constant.
Consequently there are no unmarked
multiple fibres.

\par\medskip\noindent\textbf{The exact topology and the one period shared by all leaves.}
The primary and secondary inverse-transform
charts (C.SN.3)--(C.SN.6) apply to both
contractions in (C.AS.2): they use their
identical axis normal character. For a
nonlinear leaf the explicit deformation
of its resonant coefficient to zero gives
the relative smooth-pair comparison
proved in C.2.UN.49. For a diagonal leaf
the product complement is already exact.
Thus every Hopf complement kills the
same displayed complementary meridian,
and each oriented core two-torus becomes
the ordinary fibre class of the actual
inverse transform $T_0\to B$.
The nearby smoothing is its underlying
four-manifold. Its fundamental group is
the principal-bundle presentation
(C.NS.13), with the unaltered integer
Chern vector. A two-torus has abelian
fundamental group; the injected fibre
lattice in that presentation therefore
forces (C.AS.4).

Let $\gamma$ be this one transported
integral core class. Divide the relative
canonical generator by $t^{k_P}$.
On the $i$th Hopf core its leading
term is $t^{r_i}\eta_i$, with
\[
                   \eta_i=c_i\,{dz\wedge dw\over w^{r_i+1}},
                          \qquad c_i\ne0.               \tag{C.AS.14}
\]
This is the actual coefficient calculation
of (C.NS.11), together with (C.AS.12).
Put $U=z/w^{r_i}$ and $V=\log w$.
The primitive angular period of the
secondary quotient is $(0,2\pi i/n_i)$.
For a nonlinear leaf the other period is
$(\beta_i^{-r_i},\log\beta_i)$, so
\[
       \int_{\gamma}\omega_t
          =t^{r_i}\left(\frac{\pm2\pi i c_i\beta_i^{-r_i}}{n_i}
                         +O(t)\right).                 \tag{C.AS.15}
\]
Its coefficient is nonzero. For a diagonal
leaf the other period is
$(0,\log\beta_i)$: the core lies at
constant $U$, and $c_i\,dU\wedge dV$
has zero integral there. Consequently
\[
                         \int_\gamma\omega_t=O(t^{r_i+1})
                  \quad\hbox{at a diagonal leaf}.        \tag{C.AS.16}
\]
These are periods of the same actual
class in the same nearby fibre. They
are holomorphic functions: integration
of a closed relative holomorphic form
over the transported compact cycle is
holomorphic by differentiation and
Stokes' theorem. The displayed leading
terms also give their holomorphic
extension at zero.

If a nonlinear leaf occurs, (C.AS.15)
gives one finite vanishing order $R$.
Every nonlinear $r_i$ equals it, and
(C.AS.16) gives $r_i<R$ for every
diagonal leaf. If all leaves are
diagonal, no finite vanishing order
is asserted for the given smoothing;
the explicit numerical test (C.AS.5)
is satisfied by choosing
$R=1+\max_i r_i$. This proves its
necessity without imposing an unproved
nonzero period on that family.

\par\medskip\noindent\textbf{One compact construction realizes all the specified data.}
We now use data (C.AS.1)--(C.AS.5).
The zero-Chern inverse transform is
the actual torus $T_0$ with its marked
fibres, by (C.NS.14). Take its explicit
compact $\OO_{\PP^1}(1)^{\oplus2}$
twistor family (C.NS.15), and pull it
back by $t=s^R$. For the $i$th marked
curve the exact four-period coordinate
change (C.NS.16)--(C.NS.18), with
$\gamma_i=\beta_i^{n_i}$, gives
\[
 \mathcal T_s=M_{i,s}^{\times}/
   \langle(p,\ell)\mapsto(p+d_i s^R,\gamma_i\ell)\rangle,
                    \qquad d_i\ne0 .                  \tag{C.AS.17}
\]
The non-torsion flat line and every
original primitive fibre direction
are retained by these actual period
maps.

At this marked point use weight
$(r_i,1)$ and tensor the flat line
by $\OO(n_iE_i)$, where $E_i$ is
the weighted exceptional divisor.
The same free coarse quotient as
(C.SN.9) has coordinates
$p-p_i=x^{r_i}z$, $s=xw$.
The exact infinite deck map is now
\[
 \begin{split}
 (z,w,x)\longmapsto\bigl(&
   \beta_i^{r_i}(z+d_i x^{R-r_i}w^R),\\
   &\beta_iw,\ \beta_i^{-1}x\bigr).
 \end{split}                                             \tag{C.AS.18}
\]
The exponent $R-r_i$ is nonnegative.
The added monomial has weighted degree
$r_i$, so this map commutes with the
finite action of weights
$(r_i,1,-1)$ modulo $n_i$. It preserves
$s=xw$ and changes the original base
coordinate by exactly $d_i s^R$.
Thus (C.AS.18) is the actual lift
of (C.AS.17).

On $x=0$, equality $R=r_i$ gives the
specified nonlinear contraction after
$Z=z/(\beta_i^{r_i}d_i)$. Strict
inequality $R>r_i$ gives the specified
diagonal contraction. On $w=0$ the
map is always
$(z,x)\mapsto(\beta_i^{r_i}z,\beta_i^{-1}x)$,
so the main component and its outer
trivialization are exactly
(C.SN.4)--(C.SN.5), independently of
$R$. Every given central component
and every marked normal cover is
therefore the one prescribed.

We record the analytic bounds to
verify the simultaneous compact
replacement in this varying-weight
case. On a small closed disc,
$|s|^R\le C|s|^{r_i}$ for every $i$.
The weighted norm comparison of
(C.NH.9)--(C.NH.10) consequently
gives at most
$C(1+|j|)^{n_i/r_i}$ distortion
for the $j$th translation lift.
The unitary flat metric used in
C.2.UN.50 has scalar multiplier
$|\gamma_i|\exp(-\operatorname{Re}(c_i(s)d_i s^R))$,
uniformly less than one after
shrinking the disc. Its exponential
contraction dominates that polynomial.
The same compact norm annulus proves
properness and freeness of the
infinite quotient. Coprimality in
(C.AS.1) makes its finite coarse
quotient smooth.

Outside the new Hopf component the
canonical exceptional section gives
the actual isomorphism with the
original twistor family minus its
marked central elliptic curve.
In a compact Hopf fundamental domain,
$z,w$ remain bounded and $x\to0$.
The old flat coordinate is proportional
to $x^{-n_i}$. Putting it in a compact
norm annulus uses
$O(1+|\log|x||)$ iterates, with total
base displacement
\[
                   O\bigl(|s|^R(1+|\log|x||)\bigr)
                          \longrightarrow0,
                         \qquad s=xw.                   \tag{C.AS.19}
\]
All cluster points consequently lie
on the specified original central
elliptic curve. The disjoint-curve
gluing proof of (C.NS.20) therefore
applies to all these insertions
simultaneously. It proves a smooth
Hausdorff proper local family, and
the unchanged compact twistor
complement makes the total space
compact over $\PP^1$.

The central main component is
precisely $P$, because all its
analytic torsor transitions have
been reversed by the exact inverse
charts. The diagonal automorphisms
$(z,w)\mapsto(a^{r_i}z,aw)$ commute
with either contraction and with
its finite action; they retain
every specified constant in the
axis gluing. The linear parts were
already retained by the chosen
normal-cover lattice bases.
Thus the full central support is
the specified $S$. Its equations
are $xw=s$ in the free coarse
charts, so every component has
multiplicity one. All other fibres
are the original smooth two-tori.
The Euler number of the central
support is zero by additivity:
each quasi-bundle, Hopf leaf and
double elliptic curve has Euler
number zero.

The exact symbolic checks in
\path{code/root_mixed_hopf_stars.py}
retain the weighted base map, its
Jacobian, the finite-action weights
and the common period exponent.
The finite integer search there
checks the degree exclusions in
(C.AS.9); the preceding integer
argument proves them for all weights.
\end{proof}
